\section{Typed data channels}
\label{chan-section}

We retain the labelled proof and data fragment, its fixed ambient context,
the proof-formula constructors, and the data-path wrappers constructed
above.  In particular, a formula is a stable data type sorted at the logic
vertex $b$; its proofs have mode $\mathsf p$.  All data introduced in this
section have mode $\mathsf d$.  Saying that a type is \emph{sorted at a
vertex} means that it is typed at every sort belonging to that vertex.  This
asserts the indicated typings, and does not assert that the vertex is its
exact profile.

The aim is to construct a common type of probes that works with different
choices of a data type and a data parameter. We first build observations
and predicates for one type, then query objects and a bundle that retains
the chosen parameter. These give a probe type independent of those
choices and a second observation channel on that type. The permitted
logical quantifier domains are recorded at the end of the section.

For reference, an even data path $\rho$ from a vertex at which $T$ is sorted
has the operations
\[
 \operatorname{pack}_{\rho}(t):W_{\rho}(T),\qquad
 \operatorname{use}_{\rho}(w;x.\Phi),
\]
where $t:T$, $w:W_{\rho}(T)$, and $\Phi$ is a formula under $x:T$.
An odd data path has the operations
\[
 \operatorname{mk}_{\rho}(x.\Phi):W_{\rho}(T),\qquad
 \langle w,t\rangle_{\rho}.
\]
The last expressions in both displays are formulas.  We use the proof
builders of Lemma~\ref{lem:enc-data-cancellation} in the following form:
\begin{equation}
\begin{aligned}
 \operatorname{use}_{\rho}(\operatorname{pack}_{\rho}(t);x.\Phi)
       &\leftrightarrow \Phi[x:=t] &&(\rho\text{ even}),\\
 \langle\operatorname{mk}_{\rho}(x.\Phi),t\rangle_{\rho}
       &\leftrightarrow \Phi[x:=t] &&(\rho\text{ odd}).
\end{aligned}
\label{chan-cancellation}
\end{equation}
These are logical biconditionals with typed proof terms, rather than
equations permitting replacement inside data.  Subscripts may name a
particular resulting type instead of its fixed path.  Even when paths are
reused, their source type and their role are fixed by that subscript.

\subsection{The required paths and the first channel}

Assume the configuration to be excluded: $C$ is a plain active strongly
connected component of the secondary graph, it has an internal negative
edge, and
\begin{equation}
\begin{gathered}
 (I,J,k)\text{ is a profile triple},\qquad J,k,d=\{s\}\in C,\\
 \operatorname{Ax}(s)\ne\varnothing,\qquad
 \operatorname{Ax}(s)\longrightarrow I
 \text{ by a primary positive path}.
\end{gathered}
\label{chan-configuration}
\end{equation}
The logic component $S$ has its odd closed walk and all-$S$ triple, as in
the preceding construction.  We do not assume that $C$ itself has an odd
closed walk.

\begin{lemma}[Choice of paths]
\label{chan-paths}
There is a retained negative edge associated to a triple $(p,n,g)$, with
$p,g\in C$, and there are fixed secondary paths
\[
 \alpha:k\longrightarrow p\quad\text{odd},\qquad
 \beta:p\longrightarrow k\quad\text{odd},\qquad
 \kappa:g\longrightarrow J\quad\text{even},
\]
all inside $C$.  There are also fixed secondary paths
\[
 \lambda:k\longrightarrow j\quad\text{odd},\qquad
 \theta:k\longrightarrow h\quad\text{even},
\]
where $j$ and $h$ are direct vertices.  In particular,
$\eta=\beta\theta:p\longrightarrow h$ is odd.
\end{lemma}

\begin{proof}
If $C$ has an odd closed walk, between any two of its vertices there are
paths of both parities: choose paths to and from the base of that walk,
and insert the walk once to change parity.  Choose any internal retained
negative $(p,n,g)$ and then choose the three required parities for
$\alpha$, $\beta$, and $\kappa$.

Suppose instead that every closed walk in $C$ is even.  Color its vertices
by the parity of a path from a fixed vertex.  The color is well defined:
appending a fixed return path to two such paths shows that their parities
are equal.  A positive edge preserves color and a negative edge changes
it.  The positive edge $J\longrightarrow k$ of the triple in
\eqref{chan-configuration} is retained, since $J,k\in C\subseteq H$;
hence $J$ and $k$ have the same color.  There is a negative edge entering
this color.  Indeed, take any internal negative.  If it enters the other
color, follow a directed return path to its source; that path must contain
a negative edge entering the color of $k$.  Choose such an edge
$p\longrightarrow g$, with its retained triple $(p,n,g)$.  Thus $g$ has
the color of $J,k$, and $p$ has the opposite color.  Connectivity now gives
$\kappa$ even and $\alpha,\beta$ odd.

Finally, $k$ is plain.  By the definition of plainness it has an odd path
to a direct vertex $j$ and an even path to a direct vertex $h$.  They need
not have the same endpoint and need not stay inside $C$.  Choose them as
$\lambda$ and $\theta$.  Concatenating $\beta$ with $\theta$ gives the
claimed odd path $\eta$.
\end{proof}

Fix also a path $\delta:d\longrightarrow k$ inside $C$, with no parity
condition.  We use $\delta$ only to form a type.  Work for the moment under
the data declaration $A:s$.  This is a valid declaration because
$\operatorname{Ax}(s)$ is nonempty.  Define
\begin{equation}
\begin{aligned}
 B(A)&=W_{\delta}(A),&
 L_c(A)&=W_{\alpha}(B(A)),\\
 L(A)&=W_{\lambda}(B(A)),&
 \mathsf P_L(A)&=W_{\eta}(L_c(A)).
\end{aligned}
\label{chan-basic-types}
\end{equation}
When $A$ is fixed, omit it from the notation.  The sort ledger is
\[
\begin{array}{c|c|l}
\text{data type}&\text{sorted at}&\text{support}\\ \hline
A&d=\{s\}&A:s\\
B(A)&k&\delta:d\longrightarrow k\\
L_c(A)&p&\alpha:k\longrightarrow p\text{ odd}\\
L(A)&j&\lambda:k\longrightarrow j\text{ odd}\\
\mathsf P_L(A)&h&\eta:p\longrightarrow h\text{ odd}
\end{array}
\]
Each line follows from the corresponding data-path formation lemma.
For instance, $\operatorname{mk}_{L}(t.\Phi)$ binds $t:B(A)$, whereas
$\operatorname{mk}_{\mathsf P_L}(c.\Phi)$ binds $c:L_c(A)$.  These
different sources will be important below.

\begin{definition}[Logical and raw channel values]
\label{chan-raw-log}
For $y:L$, $c:L_c$, and $P:\mathsf P_L$, put
\[
\begin{aligned}
 \operatorname{Raw}(y)
   &=\operatorname{mk}_{L_c}(t.\langle y,t\rangle_L):L_c,\\
 \operatorname{Log}(c)
   &=\operatorname{mk}_{L}(t.\langle c,t\rangle_{L_c}):L,\\
 \operatorname{Round}(y)&=\operatorname{Log}(\operatorname{Raw}(y)):L,\\
 P@y&=\langle P,\operatorname{Raw}(y)\rangle_{\mathsf P_L}.
\end{aligned}
\]
For a formula $\Phi$ under $x:L$, define
\[
 \operatorname{pred}(x.\Phi)
   =\operatorname{mk}_{\mathsf P_L}
       (c.\Phi[x:=\operatorname{Log}(c)]):\mathsf P_L.
\]
Here $t:B(A)$ and $c:L_c(A)$ in their respective callbacks.
\end{definition}

\begin{lemma}[Rounding and predicate evaluation]
\label{chan-rounding}
There are proof builders, for $y:L$ and $t:B(A)$, of
\begin{equation}
 \langle\operatorname{Round}(y),t\rangle_L
       \leftrightarrow\langle y,t\rangle_L.
\label{chan-primitive-round}
\end{equation}
For every formula $\Phi$ under $x:L$ and every $w:L$, there is a proof of
\begin{equation}
 \operatorname{pred}(x.\Phi)@w
       \leftrightarrow\Phi[x:=\operatorname{Round}(w)].
\label{chan-predicate-evaluation}
\end{equation}
\end{lemma}

\begin{proof}
For \eqref{chan-primitive-round}, cancellation for $\operatorname{Log}$
gives the comparison with
$\langle\operatorname{Raw}(y),t\rangle_{L_c}$, and cancellation for
$\operatorname{Raw}$ compares this with $\langle y,t\rangle_L$.
Apply $\mathsf{pl}$ to these two biconditionals.  For
\eqref{chan-predicate-evaluation}, expand $@$ and cancel
$\operatorname{mk}_{\mathsf P_L}$ at the exact argument
$c=\operatorname{Raw}(w)$.  Its callback becomes
$\Phi[x:=\operatorname{Log}(\operatorname{Raw}(w))]$, which is the
displayed formula by definition.
\end{proof}

\subsection{Bits stored by pairs of odd wrappers}

Two odd wrappers store one bit while retaining formula observations of
the original input. Iterating this construction will distinguish the
finite query markers.

\begin{definition}[A double wrapper]
\label{chan-double}
Let $T_1=W_{\rho}(T)$ and $T_2=W_{\sigma}(T_1)$ for two odd data paths.
Use subscripts $1$ and $2$ for their respective constructors and
evaluations.  For $t:T$ define two data embeddings
\[
 \operatorname{emb}_1(t)
   =\operatorname{mk}_2(f.\langle f,t\rangle_1),\qquad
 \operatorname{emb}_0(t)
   =\operatorname{mk}_2(f.\neg\langle f,t\rangle_1),
 \qquad f:T_1.
\]
Both have type $T_2$.  Set
\[
 \operatorname{if}(G;H,H')=(G\land H)\lor(\neg G\land H'),
\]
and, for $v:T_2$ and a formula $\Phi$ under $u:T$, set
\[
\begin{aligned}
 \operatorname{bit}(v)
   &=\langle v,\operatorname{mk}_1(u.\top)\rangle_2,\\
 \operatorname{rec}(v;u.\Phi)
   &=\operatorname{if}\bigl(\operatorname{bit}(v);
          \langle v,\operatorname{mk}_1(u.\Phi)\rangle_2,
          \neg\langle v,\operatorname{mk}_1(u.\Phi)\rangle_2\bigr).
\end{aligned}
\]
\end{definition}

\begin{lemma}[Canonical bit and recovery proofs]
\label{chan-double-laws}
For $\epsilon\in\{0,1\}$ there are proofs of the signed bit
\[
 \begin{cases}
  \operatorname{bit}(\operatorname{emb}_1(t)),&\epsilon=1,\\
  \neg\operatorname{bit}(\operatorname{emb}_0(t)),&\epsilon=0,
 \end{cases}
\]
and a proof of
\[
 \operatorname{rec}(\operatorname{emb}_{\epsilon}(t);u.\Phi)
       \leftrightarrow\Phi[u:=t].
\]
\end{lemma}

\begin{proof}
Cancel the second wrapper, then the first.  Evaluation of
$\operatorname{emb}_1(t)$ against $\operatorname{mk}_1(u.\Phi)$ is
biconditional to $\Phi[u:=t]$; evaluation of
$\operatorname{emb}_0(t)$ against the same value is biconditional to
$\neg\Phi[u:=t]$.  Each comparison follows by $\mathsf{pl}$ from the two
cancellation proofs.  With $\Phi=\top$, the proof of $\top$ and these
comparisons give the signed bits.  If $\epsilon=1$, the definition of
$\operatorname{rec}$ then selects the positive comparison.  If
$\epsilon=0$, it selects the negation of the negative comparison.  The
propositional double-negation equivalence is available from
$\mathsf{pl}$ by Lemma~\ref{lem:enc-propositional}, since the logical
constructors and classical proof builders have already been derived in
the labelled fragment.  Thus
$\mathsf{pl}$ gives the recovery biconditional in both cases.
\end{proof}

To record the iteration explicitly, begin with $T_0=L_c$ and define, for
$1\le r\le4$,
\[
 T_{2r-1}=W_{\beta}(T_{2r-2}),\qquad
 T_{2r}=W_{\alpha}(T_{2r-1}),\qquad
 \operatorname{Big}=T_8.
\]
The even-indexed types are sorted at $p$ and the odd-indexed types at $k$.
Let $\operatorname{emb}^{(r)}$, $\operatorname{bit}_r$, and
$\operatorname{rec}_r$ refer to the $r$th pair.  For a string
$\epsilon=(\epsilon_1,\ldots,\epsilon_r)$ put
\[
 E_0(c)=c,\qquad
 E_r(c;\epsilon)=\operatorname{emb}^{(r)}_{\epsilon_r}
                (E_{r-1}(c;\epsilon_1,\ldots,\epsilon_{r-1})).
\]
The bit observations at level $r$ are the formulas
\[
\begin{aligned}
 H_{r,r}(v)&=\operatorname{bit}_r(v),\\
 H_{r,a}(v)&=\operatorname{rec}_r(v;u.H_{r-1,a}(u))
                      &&(1\le a<r).
\end{aligned}
\]
Similarly, for a formula $\Phi$ under $c:L_c$, define iterated recovery by
\[
 R_0(v;c.\Phi)=\Phi[c:=v],\qquad
 R_r(v;c.\Phi)=\operatorname{rec}_r(v;u.R_{r-1}(u;c.\Phi)).
\]
Induction using Lemma~\ref{chan-double-laws} and $\mathsf{pl}$ gives the
appropriate signed proof of every
$H_{r,a}(E_r(c;\epsilon))$ and gives
\begin{equation}
 R_r(E_r(c;\epsilon);u.\Phi)
       \leftrightarrow\Phi[u:=c].
\label{chan-iterated-recovery}
\end{equation}
The induction uses logical comparisons of the displayed formulas; it
does not replace an inner encoded datum by a logically equivalent datum.

\subsection{Marked queries and point values}

Let
\[
 \mathcal T=\{G,C_0,C_\ell,C_r,E_\ell,E_r,R,\operatorname{pt}\},
 \qquad
 \mathcal T_{\mathrm{str}}=\mathcal T\setminus\{\operatorname{pt}\}.
\]
Thus there are eight marker indices and seven structural indices.
Assign distinct four-bit strings to the data tag and these markers; for
example, in the displayed order, take
\[
\begin{array}{c|ccccccccc}
\text{tag}&\mathrm{data}&G&C_0&C_\ell&C_r&E_\ell&E_r&R&\operatorname{pt}\\ \hline
\text{string}&0000&0001&0010&0011&0100&0101&0110&0111&1000.
\end{array}
\]
Write $\epsilon(q)$ for the string of a tag $q$, and put
$E_q(c)=E_4(c;\epsilon(q)):\operatorname{Big}$.  For a string
$\epsilon\in\{0,1\}^4$, let $\operatorname{bits}_{\epsilon}(v)$ be the
fixed-bracketing conjunction whose $a$th conjunct is $H_{4,a}(v)$ if
$\epsilon_a=1$ and $\neg H_{4,a}(v)$ if $\epsilon_a=0$.
Abbreviate this formula by $\operatorname{bits}_{q}(v)$ for the string of
a tag $q$.  The canonical signed-bit proofs give
$\operatorname{bits}_{q}(E_q(c))$ and
$\neg\operatorname{bits}_{q'}(E_q(c))$ whenever $q'\ne q$: in the
latter case use a position where the two strings differ and apply
$\mathsf{pl}$.

Define
\[
 \operatorname{BP}=W_{\beta}(\operatorname{Big}),\qquad
 M(A)=W_{\alpha}(\operatorname{BP}).
\]
For $P:\mathsf P_L$ and $t\in\mathcal T$, put
\begin{equation}
 \operatorname{Enc}_t(P)=\operatorname{mk}_{\operatorname{BP}}\left(
 v.\bigl(\operatorname{bits}_{\mathrm{data}}(v)
        \land R_4(v;c.\langle P,c\rangle_{\mathsf P_L})\bigr)
        \lor\operatorname{bits}_{t}(v)\right).
\label{chan-enc}
\end{equation}
Here $v:\operatorname{Big}$ and $c:L_c$, so
$\operatorname{Enc}_t(P):\operatorname{BP}$.  Choose the fixed value
\[
 c_\top=\operatorname{mk}_{L_c}(u.\top):L_c,
 \qquad u:B(A),
\]
and define, for $bp:\operatorname{BP}$,
\[
\begin{aligned}
 \operatorname{Mark}_t(bp)
    &=\langle bp,E_t(c_\top)\rangle_{\operatorname{BP}},\\
 \operatorname{Dec}(bp)
    &=\operatorname{mk}_{\mathsf P_L}
                  (c.\langle bp,E_{\mathrm{data}}(c)\rangle_{\operatorname{BP}})
                    :\mathsf P_L,\\
 \operatorname{qry}_t(o,P)
    &=\langle o,\operatorname{Enc}_t(P)\rangle_M,
                      \qquad o:M(A).
\end{aligned}
\]
The sort and value ledger for these constructions is
\[
\begin{array}{c|c|l}
\text{data type}&\text{sorted at}&\text{values just constructed}\\ \hline
T_{2r}&p&E_r(c;\epsilon)\quad(0\le r\le4)\\
T_{2r-1}&k&\operatorname{mk}\text{ values in pair }r\quad(1\le r\le4)\\
\operatorname{Big}=T_8&p&E_q(c)\\
\operatorname{BP}&k&\operatorname{Enc}_t(P)\\
\mathsf P_L&h&\operatorname{Dec}(bp)\\
M(A)&p&\operatorname{mk}_M(bp.\Phi)
\end{array}
\]
All observations, marks, queries, and recovery expressions are formulas
sorted at $b$; none is a data value of type $\operatorname{Big}$ or
$\operatorname{BP}$.

\begin{lemma}[Marker and decoder laws]
\label{chan-marker-laws}
For $P:\mathsf P_L$ and $t,u\in\mathcal T$ there are proofs of
\[
 \operatorname{Mark}_t(\operatorname{Enc}_t(P)),\qquad
 \neg\operatorname{Mark}_u(\operatorname{Enc}_t(P))\quad(u\ne t).
\]
For every $c:L_c$ there is a proof of
\[
 \langle\operatorname{Dec}(\operatorname{Enc}_t(P)),c\rangle_{\mathsf P_L}
       \leftrightarrow\langle P,c\rangle_{\mathsf P_L},
\]
and consequently, for every $y:L$, a proof of
\begin{equation}
 \operatorname{Dec}(\operatorname{Enc}_t(P))@y\leftrightarrow P@y.
\label{chan-decode}
\end{equation}
\end{lemma}

\begin{proof}
For a marker, cancel \eqref{chan-enc} at $E_u(c_\top)$.  The data-string
test is false by the signed-bit proofs.  The remaining marker-string
test holds exactly when $u=t$, again by those proofs.  Apply
$\mathsf{pl}$ to obtain the asserted signed mark.

For the decoder, first cancel its $\mathsf P_L$ constructor at the exact
input $c$, then cancel \eqref{chan-enc} at
$E_{\mathrm{data}}(c)$.  Here the data-string test holds, the
marker-string test fails, and \eqref{chan-iterated-recovery} compares the
remaining recovery formula with $\langle P,c\rangle_{\mathsf P_L}$.
Combine these proofs with $\mathsf{pl}$.  To obtain \eqref{chan-decode},
use this result at $c=\operatorname{Raw}(y)$ and expand the definition of
$@$.  Thus this calculation inserts neither
$\operatorname{Round}(y)$ nor an additional rounding of its raw value.
\end{proof}

Define the data values and the formula
\begin{equation}
\begin{aligned}
 bp_\bot&=\operatorname{mk}_{\operatorname{BP}}(v.\bot)
                       :\operatorname{BP},\\
 \operatorname{dummy}&=\operatorname{mk}_M(bp.\top):M(A),\\
 D(o)&=\langle o,bp_\bot\rangle_M,\\
 \operatorname{point}(y)
   &=\operatorname{mk}_M
       (bp.\operatorname{Mark}_{\operatorname{pt}}(bp)
                   \land(\operatorname{Dec}(bp)@y)):M(A).
\end{aligned}
\label{chan-point}
\end{equation}

\begin{lemma}[Dummy and point laws]
\label{chan-point-laws}
There are proofs of
\[
 D(\operatorname{dummy}),\qquad
 \neg\operatorname{Mark}_t(bp_\bot)\quad(t\in\mathcal T),\qquad
 \neg D(\operatorname{point}(y)),
\]
and there is a proof, for $P:\mathsf P_L$, of
\begin{equation}
 \operatorname{qry}_{\operatorname{pt}}(\operatorname{point}(y),P)
                  \leftrightarrow P@y.
\label{chan-point-query}
\end{equation}
\end{lemma}

\begin{proof}
Cancellation compares $D(\operatorname{dummy})$ with $\top$, which has
its derived proof.  It compares every marker of $bp_\bot$ with $\bot$,
giving the asserted negations by $\mathsf{pl}$.  Cancellation at
$bp_\bot$ compares $D(\operatorname{point}(y))$ with a conjunction
having $\operatorname{Mark}_{\operatorname{pt}}(bp_\bot)$ as its first
conjunct, and hence gives its negation.  Finally, cancellation compares
the left side of \eqref{chan-point-query} with
\[
 \operatorname{Mark}_{\operatorname{pt}}
       (\operatorname{Enc}_{\operatorname{pt}}(P))
 \land
 \bigl(\operatorname{Dec}(\operatorname{Enc}_{\operatorname{pt}}(P))@y\bigr).
\]
The marker proof and \eqref{chan-decode} prove the required
biconditional by $\mathsf{pl}$.
\end{proof}

\subsection{A bundle carrying an exact parameter}

Define two more odd wrappers
\begin{equation}
 F(A)=W_{\beta}(M(A)),\qquad
 \operatorname{Single}(A)=W_{\alpha}(F(A)).
\label{chan-single}
\end{equation}
Thus $F(A)$ is sorted at $k$ and $\operatorname{Single}(A)$ is sorted at
$p$.  We use $\langle f,o\rangle_F$ for $f:F(A)$ and $o:M(A)$, and
$\langle w,f\rangle_{\operatorname{Single}}$ for
$w:\operatorname{Single}(A)$ and $f:F(A)$.

\begin{definition}[Bundle and access]
\label{chan-bundle-definition}
For $m:M(A)$ and $y:L(A)$, put
\[
 \operatorname{Bundle}(m,y)=\operatorname{mk}_{\operatorname{Single}}
 \left(f.\operatorname{if}\bigl(
    \langle f,\operatorname{dummy}\rangle_F;
    \langle f,m\rangle_F,
    \langle f,\operatorname{point}(y)\rangle_F\bigr)\right)
       :\operatorname{Single}(A).
\]
Let $\Psi(o,x)$ be a formula under $o:M(A),x:L(A)$, independent of the
new variable $w:\operatorname{Single}(A)$.  Its other free data parameters
are allowed and remain fixed.  Define
\begin{equation}
\begin{aligned}
 f_i(o)&=\operatorname{mk}_{F}\left(o'.
    \operatorname{if}\bigl(D(o');\bot,
       \operatorname{qry}_{\operatorname{pt}}
           (o',\operatorname{pred}(x.\Psi(o,x)))\bigr)\right),\\
 f_e(w)&=\operatorname{mk}_{F}\left(o.
    \operatorname{if}\bigl(D(o);\top,
             \langle w,f_i(o)\rangle_{\operatorname{Single}}\bigr)\right),\\
 \operatorname{Access}(w;o,x.\Psi)
     &=\langle w,f_e(w)\rangle_{\operatorname{Single}}.
\end{aligned}
\label{chan-access-definition}
\end{equation}
Here $o,o':M(A)$, both $f_i(o)$ and $f_e(w)$ have type $F(A)$, and
$\operatorname{Access}$ is a formula.  The dependence of $f_i$ and $f_e$
on $\Psi$ is suppressed only in their names.
\end{definition}

\begin{lemma}[Access]
\label{chan-access}
Under $A:s$, $m:M(A)$, $y:L(A)$ and the proof input $d_m:\neg D(m)$,
there is a proof of
\begin{equation}
 \operatorname{Access}(\operatorname{Bundle}(m,y);o,x.\Psi)
        \leftrightarrow\Psi(m,\operatorname{Round}(y)).
\label{chan-access-law}
\end{equation}
The definitions of $\operatorname{Bundle}$ and $\operatorname{Access}$
themselves do not use $d_m$.
\end{lemma}

\begin{proof}
Throughout this proof $w$ abbreviates the exact data expression
$\operatorname{Bundle}(m,y)$; it is not a fresh datum with an assumed
property.  For any $f:F(A)$, cancellation of the bundle gives a proof of
\begin{equation}
 \langle w,f\rangle_{\operatorname{Single}}
 \leftrightarrow
 \operatorname{if}\bigl(
    \langle f,\operatorname{dummy}\rangle_F;
    \langle f,m\rangle_F,
    \langle f,\operatorname{point}(y)\rangle_F\bigr).
\label{chan-bundle-evaluation}
\end{equation}

First apply cancellation to $f_e(w)$ at
$\operatorname{dummy}:M(A)$.  It gives
\[
 \langle f_e(w),\operatorname{dummy}\rangle_F
 \leftrightarrow
 \operatorname{if}\bigl(D(\operatorname{dummy});\top,
      \langle w,f_i(\operatorname{dummy})\rangle_{\operatorname{Single}}\bigr).
\]
Together with $D(\operatorname{dummy})$ from
Lemma~\ref{chan-point-laws} and the proof of $\top$, this gives, by
$\mathsf{pl}$, a proof of
$\langle f_e(w),\operatorname{dummy}\rangle_F$.
Cancellation at the exact input $m$ gives
\[
 \langle f_e(w),m\rangle_F
 \leftrightarrow
 \operatorname{if}\bigl(D(m);\top,
                  \langle w,f_i(m)\rangle_{\operatorname{Single}}\bigr).
\]
Use $d_m$ and $\mathsf{pl}$ to obtain
\[
 \langle f_e(w),m\rangle_F
       \leftrightarrow\langle w,f_i(m)\rangle_{\operatorname{Single}}.
\]
Now use \eqref{chan-bundle-evaluation} with $f=f_e(w)$.  Its condition has
the positive proof just constructed, so $\mathsf{pl}$ selects its first
branch and gives
\begin{equation}
 \operatorname{Access}(w;o,x.\Psi)
       \leftrightarrow\langle w,f_i(m)\rangle_{\operatorname{Single}}.
\label{chan-access-first-selection}
\end{equation}
This is where the first occurrence of the callback parameter is fixed:
it is the original $m$ by the exact substitution in cancellation.

For the second selection, cancellation of $f_i(m)$ at
$\operatorname{dummy}$ gives
\[
 \langle f_i(m),\operatorname{dummy}\rangle_F
 \leftrightarrow
 \operatorname{if}\bigl(D(\operatorname{dummy});\bot,
       \operatorname{qry}_{\operatorname{pt}}
          (\operatorname{dummy},\operatorname{pred}(x.\Psi(m,x)))\bigr).
\]
The proof of $D(\operatorname{dummy})$ therefore gives a proof of
$\neg\langle f_i(m),\operatorname{dummy}\rangle_F$ by $\mathsf{pl}$.
Cancellation of $f_i(m)$ at $\operatorname{point}(y)$, followed by
$\neg D(\operatorname{point}(y))$, gives
\[
 \langle f_i(m),\operatorname{point}(y)\rangle_F
 \leftrightarrow
 \operatorname{qry}_{\operatorname{pt}}
       (\operatorname{point}(y),\operatorname{pred}(x.\Psi(m,x))).
\]
The point-query proof \eqref{chan-point-query}, at the exact predicate
$\operatorname{pred}(x.\Psi(m,x))$, compares this right side with
$\operatorname{pred}(x.\Psi(m,x))@y$.  In turn,
\eqref{chan-predicate-evaluation} compares that formula with
$\Psi(m,\operatorname{Round}(y))$.  Combining these proofs gives
\[
 \langle f_i(m),\operatorname{point}(y)\rangle_F
        \leftrightarrow\Psi(m,\operatorname{Round}(y)).
\]
Use \eqref{chan-bundle-evaluation} once more, now with $f=f_i(m)$.  Its
condition has the negative proof above, so $\mathsf{pl}$ selects its
second branch and yields
\[
 \langle w,f_i(m)\rangle_{\operatorname{Single}}
        \leftrightarrow\Psi(m,\operatorname{Round}(y)).
\]
Finally combine this with \eqref{chan-access-first-selection} by
$\mathsf{pl}$.  Only predicate evaluation rounded an argument, and that
argument was $y$.  Every occurrence of $m$ in this derivation comes from
exact data substitution, without a logical replacement of data.
\end{proof}

\subsection{The common type of probes}

Recall the retained triple $(p,n,g)$ selected in
Lemma~\ref{chan-paths}.  Its tail type $U_n$ is the previously fixed
positive data wrapper of a literal sort.  It is sorted at $n$, with the
operations $\operatorname{send}_n$ and $\operatorname{read}_n$ and the
read--send biconditional \eqref{eq:enc-read-send}.  Since
$\operatorname{Single}(A)$ is sorted at $p$, form
\begin{equation}
 \operatorname{Fun}(A)
    =\Pi^{\mathsf d,\mathsf d}w:\operatorname{Single}(A).U_n,
 \qquad\text{sorted at }g.
\label{chan-fun}
\end{equation}
This uses precisely the formation rules represented by $(p,n,g)$.

Let $\zeta$ be the fixed primary positive path
$\operatorname{Ax}(s)\longrightarrow I$ of
\eqref{chan-configuration}, and positively wrap the literal $s$:
\begin{equation}
 Z=W_{\zeta}^{+}(s),\qquad
 A^\uparrow:Z\quad(A:s),\qquad
 z^\downarrow:s\quad(z:Z),\qquad
 (A^\uparrow)^\downarrow={}_\beta A.
\label{chan-z}
\end{equation}
The initial type $s$ is sorted at all of $\operatorname{Ax}(s)$, so $Z$
is sorted at $I$.  The superscripts in \eqref{chan-z} denote the fixed
positive lift and projection, not additional primitive operations.

Now define the data type
\begin{equation}
 V=\Pi^{\mathsf d,\mathsf d}z:Z.
             W_{\kappa}(\operatorname{Fun}(z^\downarrow)).
\label{chan-v}
\end{equation}
Under $z:Z$ the projection $z^\downarrow:s$ can be substituted for the
data parameter $A$ in every preceding channel type.  The resulting
$\operatorname{Fun}(z^\downarrow)$ is sorted at $g$, and its wrapper
along $\kappa:g\longrightarrow J$ is sorted at $J$.  The triple
$(I,J,k)$ therefore sorts $V$ at $k$.  This construction binds the former
parameter $A$ through $z^\downarrow$: $V$ has no free $A$ or $m$.

For $v:V$, $A:s$, $m:M(A)$ and $y:L(A)$, define the formula
\begin{equation}
 v\diamond_{A,m}y
   =\operatorname{use}_{\kappa}\left(
        v\,A^\uparrow;
        f.\operatorname{read}_n(f\,\operatorname{Bundle}(m,y))\right).
\label{chan-diamond}
\end{equation}
The applications displayed here have label $(\mathsf d,\mathsf d)$.
For clarity, their typing chain is
\[
\begin{aligned}
 v\,A^\uparrow
   &:W_{\kappa}(\operatorname{Fun}((A^\uparrow)^\downarrow))
       =_\beta W_{\kappa}(\operatorname{Fun}(A)),\\
 f&:\operatorname{Fun}(A),\qquad
 \operatorname{Bundle}(m,y):\operatorname{Single}(A),\\
 f\,\operatorname{Bundle}(m,y)&:U_n.
\end{aligned}
\]
The conversion in the first line is allowed because all paths, tails,
and role-dependent syntax choices were fixed before substitution.  Since
$\kappa$ is even, its $\operatorname{use}$ operation accepts exactly the
formula callback in \eqref{chan-diamond}.  Thus $\diamond$ is well typed
without a proof assumption about $m$.

The final data-sort ledger is
\[
\begin{array}{c|c|l}
\text{data type}&\text{sorted at}&\text{formation support}\\ \hline
M(A)&p&W_{\alpha}(\operatorname{BP}(A))\\
F(A)&k&W_{\beta}(M(A))\\
\operatorname{Single}(A)&p&W_{\alpha}(F(A))\\
U_n&n&\text{the retained negative's fixed tail}\\
\operatorname{Fun}(A)&g&(p,n,g)\\
Z&I&\zeta:\operatorname{Ax}(s)\longrightarrow I\text{ positive}\\
W_{\kappa}(\operatorname{Fun}(z^\downarrow))&J&\kappa:g\longrightarrow J\text{ even}\\
V&k&(I,J,k)
\end{array}
\]
In particular, neither the use of $M(A)$ as an open parameter nor that of
$V$ as an input type requires $p$ or $k$ to be direct.

\subsection{The second channel and permitted quantifiers}

Apply the first-channel construction to $V$ itself, which is sorted at
$k$, using the same three fixed odd paths:
\begin{equation}
 X_c=W_{\alpha}(V),\qquad
 X=W_{\lambda}(V),\qquad
 \mathsf P_X=W_{\eta}(X_c).
\label{chan-x-types}
\end{equation}
These types are sorted at $p,j,h$, respectively.  For $x:X$ and $c:X_c$
define
\[
\begin{aligned}
 \operatorname{Raw}_X(x)&=\operatorname{mk}_{X_c}
                          (v.\langle x,v\rangle_X),\\
 \operatorname{Log}_X(c)&=\operatorname{mk}_{X}
                          (v.\langle c,v\rangle_{X_c}),\\
 \operatorname{Round}_X(x)&=\operatorname{Log}_X(\operatorname{Raw}_X(x)),\\
 i@x&=\langle i,\operatorname{Raw}_X(x)\rangle_{\mathsf P_X}
                         \qquad(i:\mathsf P_X),\\
 \operatorname{pred}_X(x.\Phi)
   &=\operatorname{mk}_{\mathsf P_X}
                          (c.\Phi[x:=\operatorname{Log}_X(c)]).
\end{aligned}
\]
The callbacks defining $\operatorname{Raw}_X$ and
$\operatorname{Log}_X$ bind $v:V$.  Write
\[
 [xv]=\langle x,v\rangle_X\qquad(x:X,\ v:V).
\]
The proofs of Lemma~\ref{chan-rounding}, with these types substituted,
give
\begin{equation}
\begin{aligned}
 {}[\operatorname{Round}_X(x)v]&\leftrightarrow[xv],\\
 \operatorname{pred}_X(x.\Phi)@w
       &\leftrightarrow\Phi[x:=\operatorname{Round}_X(w)].
\end{aligned}
\label{chan-x-evaluation}
\end{equation}

All logical quantifiers in the subsequent relational construction have
one of the following fixed domain annotations:
\[
\begin{array}{c|c|c}
\text{domain}&\text{chosen sorting vertex}&\text{quantifier constructor}\\ \hline
L(A)&j&\text{the fixed direct-domain constructor at }j\\
X&j&\text{the fixed direct-domain constructor at }j\\
\mathsf P_L(A)&h&\text{the fixed direct-domain constructor at }h\\
\mathsf P_X&h&\text{the fixed direct-domain constructor at }h.
\end{array}
\]
Here $j$ and $h$ are direct by Lemma~\ref{chan-paths}, so each constructor
has its required profile triple and fixed proof paths.  Existentials use
the same domain annotation through their derived definition.  Other
generic data, including $A:s$, $m:M(A)$, and $v:V$, occur as open
parameters or as raw data binders whose supporting products have been
displayed above.  Schematic assertions about these parameters denote
proof builders in valid contexts, not further logical quantifiers.
Every data definition in this section is independent of proof variables;
the proof input $d_m$ is used only in the proof of
Lemma~\ref{chan-access}.
